% AUTO-GENERATED by paper/analysis/make_prompt_appendix.py -- do not hand-edit.
% Rendered from the live prompt builders so the appendix cannot drift from the
% implementation it discloses. Regenerate with:
%   python paper/analysis/make_prompt_appendix.py

The three prompts below are reproduced as the system emits them, with only the
per-run document content replaced by a placeholder. Together with the schema of
Appendix~\ref{sec:appendix-schema} they are what this paper discloses in place of
a code release (\S\ref{sec:artifacts-statement}).

\subsection{Pass 1: depth-map inference}
\label{sec:prompt-depthmap}
\begin{lstlisting}[style=promptstyle]
You are analyzing the structural skeleton of an early learning standards document. You will NOT extract individual elements — only the document's nesting hierarchy.

Read the sample below and identify how many distinct nesting depths the document uses, what each depth looks like (its prefix/label pattern), and one concrete example. The DEEPEST depth is always the leaf (individual learning goals/foundations/benchmarks).

This document's depths must then be mapped to a canonical 4-level hierarchy:
- depth 1 → "domain"
- depth 2 → "strand"
- depth 3 → "sub_strand"
- deepest depth → "indicator"

Rules:
- Fill the canonical levels FROM THE TOP DOWN. The first depth is always "domain", the deepest depth is always "indicator", and any depth in between takes the next canonical level in order: the depth directly under domain is "strand", the one below that is "sub_strand". A document with fewer depths therefore drops the levels at the BOTTOM of that middle range — never "strand".
- A 4-level document (domain > group > sub-group > leaf) maps to: domain, strand, sub_strand, indicator.
- A 3-level document (domain > group > leaf) maps to: domain, strand, indicator. There is NO sub_strand — the single middle level is a STRAND, never a sub_strand.
- A 2-level document (domain > leaf) maps to: domain, indicator. There is no strand and no sub_strand.
- Use depth POSITION, not the document's labels. If a document uses "Sub-Strand" as the label for the second level (directly under domain), it is still a STRAND in our canonical hierarchy — and likewise a document that calls its single middle level "Section", "Concept", "Goal", or anything else still maps that level to "strand".

Every line is prefixed with its page and the horizontal position it was laid out at: `[Page N | x=0.09]`, where x is the line's left edge as a fraction of the page width. Use x to work out whether a page uses a side-by-side layout: lines of one column share nearly the same x, and a jump to a clearly different x means a different column. Lines you judge to be in the same column read as one continuous cell, even when lines from other columns are interleaved between them. Treat parallel columns as repetitions of ONE depth, not as separate depths.

Be deterministic and conservative. Do not speculate. If you cannot tell whether two depths are distinct, assume they are the same depth.

Output ONLY a JSON object with this exact shape:
{
  "doc_depths": [
    {"depth": 1, "canonical_level": "domain",     "label_in_doc": "<what the doc calls this>", "prefix_pattern": "<regex-ish pattern e.g. 'ALL-CAPS HEADING' or 'N. <Title>: <desc>'>", "example": "<exact text from doc>"},
    {"depth": 2, "canonical_level": "strand|indicator", ...}
  ],
  "notes": "<one short sentence about the LEVEL LAYOUT only, e.g. 'age-band columns repeat one depth'. A later stage reads this as an instruction, so: never describe text below the deepest depth as examples/illustrations/anecdotes, and never say to ignore or skip anything. Empty string if nothing unusual.>"
}

DOCUMENT SAMPLE:

<<serialized document lines, one per row, each prefixed
  [Page N | x=<left edge as fraction of page width>]>>
\end{lstlisting}

\subsection{Pass 2: per-chunk detection}
\label{sec:prompt-detection}
\begin{lstlisting}[style=promptstyle]
You extract structural elements from an early learning standards document chunk.

Be deterministic. Be conservative. Do not be creative. Two runs over the same text MUST produce the same JSON. Do not invent titles, codes, or descriptions — only use text that literally appears in the chunk.

CANONICAL HIERARCHY: domain (1) > strand (2) > sub_strand (3) > indicator (leaf). A document with fewer levels drops them from the BOTTOM of the middle: the level directly under domain is always `strand`, and `sub_strand` exists only when there are two levels between domain and the leaf. So a 3-level document is domain > strand > indicator, never domain > sub_strand > indicator. The depth map says which levels exist.

LAYOUT COORDINATES: every line is prefixed with its page and the horizontal position it was laid out at — `[Page N | x=0.09]`, where x is the line's left edge as a fraction of the page width (0.00 is the left margin, 1.00 the right). Use x to work out the page's column structure yourself: on a side-by-side layout the lines of one column share nearly the same x, and a jump to a clearly different x means a different column. Lines you judge to be in the same column must be read together as one continuous cell, in the order given, even when lines from other columns are interleaved between them. Never join text across what you judge to be two different columns into a single element, and never attribute one column's prose to another column's element. Expect a header that spans the whole layout to sit at the leftmost x while belonging to every column beneath it. A line tagged only `[Page N]` has no usable coordinate; use the text's own layout cues there.

CLASSIFICATION RULE:
- If a depth map is provided, look up each element's depth in the document and use the `canonical_level` from that depth. Do NOT reclassify based on prefix style.
- If no depth map is provided, classify by nesting POSITION, never by what the document calls a level.

EXTRACTION RULES:
1. Emit every structural element you see, even if its children are not in this chunk. Emit ONLY what you see: every element must be backed by a heading line that literally appears in this chunk, and its `title` must be text copied from that line. A position identifier you read inside a CHILD's code is NOT evidence that the parent's heading is present. A child code names its ancestors — an indicator coded `A.B.2.c` says it sits under the `2` at the level above — but that tells you only where the child belongs; it does not put the parent's heading on the page. When a chunk opens part-way through a section and an ancestor's heading was left behind in the previous chunk, emit the children alone and leave the ancestor out. A later chunk that contains the heading will supply it, and the parent is reconstructed downstream from the children's codes. Never reconstruct it here by reading the code apart. Emit ONLY what you see: every element must be backed by a heading line that literally appears in this chunk. (This says only that the heading must BE there — how you split that line into `title`, `code`, and `description` is decided by rules 4 and 8 as usual, and this rule never tells you to copy the whole line into `title`.) A position identifier you read inside a CHILD's code is NOT evidence that the parent's heading is present. A child code names its ancestors — an indicator coded `A.B.2.c` says it sits under the `2` at the level above — but that tells you only where the child belongs; it does not put the parent's heading on the page. When a chunk opens part-way through a section and an ancestor's heading was left behind in the previous chunk, emit the children alone and leave the ancestor out. A later chunk that contains the heading will supply it, and the parent is reconstructed downstream from the children's codes. Never reconstruct it here by reading the code apart.
2. A lettered/bulleted list (a., b., c., …) is EITHER the leaf indicators OR illustrative examples below a leaf — decide using the depth map, never the "a./b." surface style:
   - If the document's DEEPEST depth in the depth map is this lettered list (i.e. there is no deeper indicator level beneath the letters), then EACH lettered item IS an indicator. Emit one indicator per letter. Its `title` is the lettered item's own skill/competency statement (the text after the letter, e.g. "a. Demonstrates self-confidence." → "Demonstrates self-confidence"); any further example anecdotes indented beneath that letter go in its `description`. Do NOT fold these letters away and do NOT drop them — every lettered item in the chunk must appear.
   - Otherwise (the letters sit BELOW an already-identified leaf indicator and read as concrete behavior anecdotes, not skill statements) they are an EXAMPLE LIST — see rule 2a. Do NOT emit them as separate indicators.
2a. EXAMPLE LISTS — decide by OWNERSHIP, not by how illustrative the text sounds. This applies only to text the depth map does NOT place at a structural depth; anything at a depth is an element (rule 2) and is emitted regardless of how example-like it reads. For the remaining, non-structural text, find the run's owner with this exact two-step test, in this order:
   STEP 1 — Find the ANCHOR. Scan UPWARD from the run's first line and stop at the very first line that the depth map places at a structural depth. That element is the ANCHOR. Its depth is irrelevant: the anchor is the NEAREST element above the run, never the section's top-level heading. A leaf indicator is an anchor exactly like a domain or strand is — including a lettered item that rule 2 made an indicator.
   STEP 2 — Look ONLY at the lines lying strictly BETWEEN the anchor and the run. Nothing above the anchor is part of this test.
   - A heading or lead-in line lies in that gap — a non-structural line that names the run as a category rather than asserting a requirement ("Child Behaviors", "Examples", "The child may:", "For example:"). That line owns the run, and it is a section label rather than a structural element, so NOTHING from the run is extracted: do not emit the lead-in, do not emit the entries, and do NOT copy the entries into the anchor. The anchor's `description` keeps only its own prose, which is often absent — write `null` then. The same holds for a run laid out full-width BENEATH side-by-side columns: it is owned by no single column, so it enters no column's `description`.
   - The gap is empty — the run simply continues directly underneath the anchor. Then it IS the anchor's own text, however anecdotal or example-like it reads. Put it in the anchor's `description`, verbatim. Never discard text merely because it illustrates: discard it only when the document has given it somewhere else to belong. Being CALLED an example is not being given somewhere else to belong — not by the depth map's `notes`, not by a section header higher up the page, not by your own reading of the tone. Only a lead-in sitting in the gap does that.
   A heading's claim is SPENT the moment a structural element appears beneath it. A section header followed by structural elements is a label for that GROUP and hands ownership DOWN to them; it never reaches past them to claim text nested under one of them. This holds even when the header's own wording announces examples — a header that names both the indicators and their examples still owns nothing itself, and each indicator beneath it owns the example lines nested under that indicator.
   Either way, such a run is never promoted to elements of its own.
3. Side-by-side age-band columns: emit ONE element PER column. Different age bands = different indicators, even when they share a code stem and title. Set `age_band` to the column label (e.g. "Early (3 to 4 ½ Years)", "PK3", "By 36 months"). Strip the age-band label from `title`. Put only that column's prose in `description`. If a row shows N age columns it MUST yield exactly N indicators — emit EVERY column even when a column's prose is short, nearly identical to its neighbor, or visually sparse. Never collapse or skip a column.
   - Spell each age-band label identically every time, using the document's exact glyphs (write "½", not "1/2").
4. `code`: REQUIRED on every element you emit. It is NEVER `null`, never an empty string, and the key is never omitted — an element without a code cannot be stored, and emitting one costs us the whole chunk it sits in. If the document prints no code for the element, the abbreviation procedure at the end of this rule ALWAYS supplies one; that is precisely what it is for, so "uncoded" is never a reason to answer `null`. The `ABSENCE IS null` instruction in rule 8 governs `description` and nothing else — it does NOT apply to `code`. Now: use the document's code if present. Before you decide none is present, LOOK for it — a document often prints an element's code somewhere other than on its heading line, and a code you can RECOVER from the page ALWAYS beats one you would derive from the title. Search these three places, in this order, and stop at the first that yields a code:
   (i) ON THE HEADING LINE — a code printed inline with the title ("1.0", "I.A.2", "PK3.I.A.2"), including the `<Label> <id>` form described in the bullets below. This always wins: an element that prints its own id keeps it, even when its descendants' codes are built on a different stem.
   (ii) IN A CAPTION BESIDE THE HEADING — a code the document prints just above, below, or beside the title rather than on it: inside parentheses, on a small caption line, in a table or column header, or in a lead-in that names the group of items which follows. A heading "Working Memory" whose next line reads "Indicators (CD.WM)" has code "CD.WM". Take only the IDENTIFIER — the caption's surrounding words ("Indicators", "Standards", "Codes") name what the caption points at and are not part of the code — and copy it VERBATIM, dots included.
   (iii) FROM ITS DESCENDANTS' CODES — the heading prints no code of its own anywhere, but the elements BENEATH it carry dotted codes that all begin with the same segments. An ancestor's code is the COMMON PREFIX of its descendants' document codes. Indicators coded "CD.WM.PK1", "CD.WM.PK2", "CD.AT.PK1", "CD.AT.PK2" say that the group heading above the first two is "CD.WM", the group above the last two is "CD.AT", and the level above all four is "CD". Emit the WHOLE shared prefix with its dots ("CD.WM") — never just its last segment ("WM"), and never an abbreviation of the title.
       Peel exactly ONE segment per level as you move UP, so a recovered prefix never reaches past the level directly below it: the parent of an indicator coded "CD.WM.PK1" is "CD.WM", and the level above THAT is "CD". A level that prints its own code under (i) takes that code and consumes NO segment — keep peeling from where you were for the level above it. Peel this way even when the chunk happens to show only one group under a heading, where the raw shared prefix would over-reach ("CD.WM" for a level that is really "CD").
       Take a prefix only when the evidence is real:
       - The descendant's code must have MORE THAN ONE dot-separated segment. A single-segment code ("1", "a", "VOCA") names only the element itself and holds no ancestor inside it.
       - At least TWO descendants must agree on the prefix AND differ after it. One descendant alone cannot show where its code stops naming ancestors and starts naming itself.
       - Only WHOLE segments count. Never split a segment.
       - A descendant code that carries a structural LABEL word ("Foundation 1.7", "Benchmark 1.1", "Standard 2") is NOT a namespace path: the label names that descendant's OWN level, so the entire id after the label belongs to the descendant and none of it is an ancestor's code. Take a prefix only from bare dotted codes ("CD.WM.PK1", "I.A.2").
       This rule decides WHICH code an element gets. It NEVER licenses emitting an element you cannot see — rule 1 still governs that, and the heading line must literally appear in this chunk. If it does not, emit the descendants alone and take nothing from them.
   When (ii) or (iii) supplied the code, cite the line it came from in `source_text` IN ADDITION to the heading line — that line is part of the text you used (rule 7). Never instead of the heading line: rule 1's self-check still holds, and if a descendant's line is the ONLY line you can cite, you did not see the heading and must not emit the element at all.
   ONLY when all three come up empty — the document leaves this element genuinely uncoded — do you generate a stable ≤5-char uppercase abbreviation from the title, by this EXACT procedure:
   (a) Split the title into words on spaces and slashes. A hyphenated compound is ONE word ("self-selected" → one word contributing "S"; "Health/Mental Wellness" → three words).
   (b) DROP every CONNECTOR word. A connector is any of: `a an the and or but nor of to in on at by for from with into about over under through as`, and the symbol `&`. Connectors carry no meaning in a code — their initials only crowd out the words that do, and the code has just 5 characters to spend. Keeping them yields a code made of function words: "Attends to an adult or peer who is communicating verbally or nonverbally" would give "ATAAO" ("Attends To An Adult Or") instead of "AAPWI", and "Engages in an activity for a sustained period of time" would give "EIAAF" instead of "EASPT". The code is the human-readable part of a standard's identifier, so it must spend its characters on the words that name the skill. This list is exact — do not extend it to other prepositions you may consider connector-like ("toward", "during", "between", "despite", "while"): those are CONTENT words for this purpose and their initials stay in the code. Match case-insensitively and ignore trailing punctuation.
   (c) If exactly ONE content word remains, the code is its first 4 letters ("Vocabulary" → "VOCA", "Initiative" → "INIT"). Otherwise the code is the first letter of each remaining content word, in order, capped at the first 5 ("Physical Development" → "PD", "Approaches to Learning" → "AL", "Concepts About Print" → "CP", "Social & Emotional Development" → "SED", "Language and Early Literacy" → "LEL", "Engages in an activity for a sustained period of time" → "EASPT").
   (d) If dropping connectors would leave NO words at all, keep the title's words as they are and apply (c) to those.
   Use the SAME code every time the same element appears.
   - When a heading is written as `<Label> <id>: <Title>` — where `<Label>` is ANY structural word the document uses to name a level (e.g. Strand, Concept, Sub-Strand, Goal, Pillar, Unit, Theme, Section, Element, Standard, Domain, Benchmark, Component, Cluster, Foundation, Outcome, or whatever else this particular document uses) and `<id>` is the position identifier at that level (a number "1", a letter "A", a token "1.2") — the label-plus-id together is the element's `code`, written as the label, ONE space, then the id (e.g. "Strand 1", "Concept 2", "Goal A", "Pillar 1.2"). Whatever punctuation the heading used to separate the label from the id — a colon, hyphen, en/em dash — is layout and does NOT enter the code: "Strand: 1.0 — Listening and Speaking" gives code "Strand 1.0", never "Strand: 1.0". The `title` is ONLY the text after the colon (e.g. "Strand 1: Self-Awareness and Emotional Skills" → title "Self-Awareness and Emotional Skills"). This rule applies to EVERY structural-label heading word, not just the examples — never leave the `<Label> <id>:` prefix inside `title`, regardless of which word the document chose for `<Label>`.
   - A structural label word contributes NOTHING to the code on its own. It earns a place in the code only through the position identifier attached to it: "Strand 1" works as a code because "1" says WHICH strand. So when a heading names the level but supplies no position identifier — `<Label><separator><Title>` with nothing between the label and the title, whatever the separator (colon, hyphen, en/em dash, or just a space): e.g. "Sub-Strand — Vocabulary", "Concept - Curiosity and Interest", "Domain: Number Sense", "Goal — Working Memory" — the label is pure level-naming, exactly like the trailing structural noun below. Drop it and derive the code from the TITLE ALONE using the ≤5-char abbreviation rule ("Vocabulary" → "VOCA", "Curiosity and Interest" → "CI", "Number Sense" → "NS", "Working Memory" → "WM") — unless the document supplies a code for the element somewhere else, in a caption or through its descendants' codes, in which case that recovered code wins over any abbreviation ((ii)/(iii) above). Never abbreviate the label into the code and never prefix the code with it ("Sub-Strand — Vocabulary" → "VOCA", NOT "SS-V"; "Concept - Curiosity and Interest" → "CI", NOT "C-CI"), and never use the raw heading line as the code.
   - The same principle holds when the structural label appears as a TRAILING noun on the heading instead of a prefix: a heading like `<Title> <Label>` (e.g. "Social and Emotional Development Domain", "Language and Literacy Standard", "Physical Development Area", "Number Sense Strand") names a `<Title>` of a level the document calls `<Label>`. The structural noun is the level word, NOT part of the name — emit only the bare `<Title>` ("Social and Emotional Development", "Language and Literacy", "Physical Development", "Number Sense"). This applies regardless of which structural noun the document uses (Domain, Standard, Area, Strand, Section, Component, Cluster, Foundation, etc.). The same-named element may also appear elsewhere on the page as a running page header in ALL CAPS (e.g. "SOCIAL EMOTIONAL DEVELOPMENT STANDARD" repeated at the top of every page) — that running header is page typography, not a separate element and not a code; if you emit anything for it, normalize it back to the same bare title-cased `<Title>` so the two variants share the SAME `code` and the SAME `title`. The abbreviation rule above ("≤5-char uppercase abbreviation from the title") still applies WHEN the document leaves the element uncoded: derive the code from the BARE title (e.g. "Social Emotional Development" → "SED"), NEVER use the heading text itself or the structural noun as the code. If the document DOES code it — in a caption, or through the shared prefix of its descendants' codes — that recovered code wins ((ii)/(iii) above).
   - When a lettered list IS the leaf indicators (per rule 2), the `code` is JUST that item's letter, lowercased ("a", "b", "c", …) — exactly as the document orders them, regardless of any OCR casing. Do NOT prepend the parent strand/concept number (no "S1C1a"), and do NOT uppercase it (no bare "C"). The downstream parser supplies the parent's number; the detector only needs the consistent local letter.
5. `confidence`: 0.95+ if the depth map clearly applies; 0.80-0.94 if the chunk is ambiguous but the answer is likely; <0.70 if you are guessing.
6. `source_page`: page number from the `[Page N]` / `[Page N | x=…]` marker on that line.
7. `source_text`: the exact line(s) from the chunk you used. Copy verbatim, WITHOUT the leading `[Page N]` / `[Page N | x=…]` marker.
8. `description`: capture the element's COMPLETE descriptive/introductory prose, verbatim — EVERY sentence of a domain/strand/sub_strand introduction that appears in the chunk, not just the first sentence. Do NOT summarize, paraphrase, shorten, or stop at the first sentence. If the intro runs across several sentences or lines in the chunk, include all of them. (For an age-band column indicator, `description` is still only that one column's prose, per rule 3.) `description` holds only the text the element OWNS (rule 2a) — never a run that a heading or lead-in of its own has claimed; an element that owns no prose gets `null` for `description`, and that is the correct answer, not a gap to fill. Conversely, do not null out a `description` just because the text it owns reads like examples — unclaimed text under an element is that element's own.
   - ABSENCE IS `null`, NEVER `""`. When an element owns no prose, write `null` — not an empty string, not a string of spaces. `""` and `null` are two spellings of the same fact, and emitting both across a document makes the same absence irreconcilable downstream. Use `null` every time.

NEGATIVE EXAMPLES (do NOT do these):
- Do not emit "Indicators and Examples in the Context of Daily Routines" as a structural element. It is a section header. It is not an owner either: structural indicators follow it, so its claim is spent on them (rule 2a) — the example lines nested under each of those indicators belong to that indicator, not to this header.
- Do not merge "Early" and "Later" age columns into one indicator.
- Do not emit `"code": null` (or `""`, or omit the key) for an element the document leaves uncoded. Derive the ≤5-char abbreviation from the title instead: "Nutrition" → "NUTR", not `null`. A null code is never the right answer for any element at any level.
- Do not keep a structural label inside the title: "Strand 1: Self-Awareness" → title is "Self-Awareness", NOT "Strand 1: Self-Awareness".
- Do not keep a trailing structural noun inside the title: "Social and Emotional Development Domain" → title is "Social and Emotional Development", NOT "Social and Emotional Development Domain". "Language and Literacy Standard" → title is "Language and Literacy".
- Do not classify a numeric prefix ("1.", "2.") as `sub_strand` just because numeric-under-letter is sub_strand in some other doc — use the depth map.
- Do not truncate a multi-sentence domain/strand description to its first sentence — capture the entire introduction verbatim.
- When the depth map's leaf is a lettered list, do NOT drop or fold its letters: every "a./b./c." skill statement under a concept is its own indicator (e.g. under "Phonological Awareness" emit indicators "a","b","c","d","e","f","g", one per letter). Emit them even for concepts whose letters appear far from the concept heading in the chunk.
- When a lettered list is NOT the leaf — concrete anecdotes ("Chooses carrots over celery during mealtime.") sitting under an already-identified leaf indicator/foundation — do NOT promote them to indicators. Whether they land in that indicator's `description` is decided by rule 2a's ownership test, not by how anecdotal they sound.
- Do not drop text that no heading claimed. Anecdote lines running on directly beneath a leaf ("Acknowledges her own accomplishments and says, "I can hit the ball."") — with no "Examples:"-style lead-in between them and the leaf — are that leaf's own `description`. A section header further UP the page does not change this: the leaf is the nearest element above the anecdotes, the gap between them is empty, so they are the leaf's. Emitting `null` there loses real content. Apply this identically to every such leaf on every page — two pages with the same layout must get the same answer.
- Do not paste a claimed run into the element above it. Under a "Child Behaviors" heading and a "The child may:" lead-in, the behaviors belong to that section — they go nowhere, and the indicator's `description` is `null`.
- Do not treat a phrase that merely announces an upcoming example list as an element or as a description. When the list follows that phrase directly, the phrase introduces illustrations and goes wherever the list goes: nowhere. When a structural element intervenes between the phrase and the list, the phrase's claim is already spent — the list is that element's own text (rule 2a).
- Do not fold a structural label into the code: "Sub-Strand — Vocabulary" → code "VOCA", never "SS-V", "SS-VOCA", or the whole heading line "Sub-Strand — Vocabulary". A label with no position identifier after it is not part of the code. Give the SAME element the SAME code on every chunk it appears in — never "VOCA" in one chunk and "SS-V" in another.
- Do not invent an abbreviation for an element the document has already coded elsewhere on the page. A group heading "Working Memory" captioned "Indicators (CD.WM)", or one whose indicators all read "CD.WM.PK1", "CD.WM.PK2", has code "CD.WM" — not "WM" (its last segment alone), and not "WORK" or "WM" derived from its title. The abbreviation is the LAST resort, for an element the document leaves genuinely uncoded; a code printed anywhere that points at the element beats it every time. The same holds for the level ABOVE that group: a domain heading with no printed code whose indicators read "CD.WM.PK1" and "CD.AT.PK1" has code "CD", not an abbreviation of its title.
- Do not let a recovered prefix reach past the level directly below it. If the chunk shows a domain heading and only ONE group beneath it, the indicators' shared prefix "CD.WM" is the GROUP's code, not the domain's — peel one segment per level and give the domain "CD". A prefix that equals the code of the element directly below it is a level too long.
- Do not read an ancestor's code out of a descendant's LABELLED id. Indicators coded "Foundation 1.7" and "Foundation 1.8" under a sub_strand "Grammar" do NOT give that sub_strand the code "Foundation 1": the label "Foundation" names the INDICATOR's level, so "1.7" is the indicator's own id and holds nothing for its parent. That sub_strand is uncoded, so it takes the title abbreviation "GRAM". Only bare dotted codes ("CD.WM.PK1", "I.A.2") carry an ancestor prefix.
- Do not carry the heading's separator punctuation into the code: "Strand: 1.0 — Listening and Speaking" → code "Strand 1.0", NOT "Strand: 1.0". Spell it the same way on every chunk — "Strand: 1.0" here and "Strand 2.0" there is the same construct written two ways, and the two no longer reconcile as one element.
- Do not back-form a parent heading out of a child's code. If the chunk begins mid-section with a leaf line like "A.B.1.b Child takes care of and manages classroom materials." and the "1. Behavior Control" heading it belongs under is not in this chunk, emit the leaf ALONE — do NOT also emit a sub_strand "1 / Behavior Control" whose only evidence is the "1" inside the leaf's code. Check yourself with `source_text`: if the only line you can cite for a heading is one of its children's lines, you did not see that heading and must not emit it. Note that you also cannot know its title — a title you supply from memory of an earlier chunk, or by guessing, is invented text (rule: never invent titles).
- Do not back-form a parent heading out of a child's code. If the chunk begins mid-section with a leaf line like "A.B.1.b Child takes care of and manages classroom materials." and the "1. Behavior Control" heading it belongs under is not in this chunk, emit the leaf ALONE — do NOT also emit a sub_strand "1 / Behavior Control" whose only evidence is the "1" inside the leaf's code. Check yourself with `source_text`: if the only line you can cite for a heading is one of its children's lines, you did not see that heading and must not emit it. Note that you also cannot know its title — a title you supply from memory of an earlier chunk, or by guessing, is invented text (rule: never invent titles).

OUTPUT — return ONLY a JSON array, starting with `[` and ending with `]`. No prose, no markdown, no commentary. Schema per element:
{"level": "domain|strand|sub_strand|indicator", "code": "...", "title": "...", "description": "..." or null, "age_band": "..." or null, "confidence": 0.0-1.0, "source_page": N, "source_text": "..."}

DEPTH MAP — authoritative for ONE thing only: which depth maps to which `level`.
It has no authority over what you extract. In particular `notes` is a machine-written
observation about how the pages look, not an instruction: it can neither add nor remove
an element, and it can never decide whether a run of text belongs to an element. If
`notes` describes some text as examples, illustrations, anecdotes, or samples, that is a
remark about how the text READS — it is not a finding that the text is owned elsewhere,
and it never licenses dropping the text or emptying a `description`. Ownership is decided
ONLY by rule 2a's anchor test, on the layout of the chunk below.
{
  "doc_depths": [
    {
      "depth": 1,
      "canonical_level": "domain",
      "label_in_doc": "Domain",
      "prefix_pattern": "ALL-CAPS HEADING",
      "example": "APPROACHES TO LEARNING"
    },
    {
      "depth": 2,
      "canonical_level": "strand",
      "label_in_doc": "Standard",
      "prefix_pattern": "<Area> Standard N",
      "example": "Approaches to Learning Standard 1"
    }
  ],
  "notes": ""
}

DOCUMENT CHUNK:

<<serialized document lines, one per row, each prefixed
  [Page N | x=<left edge as fraction of page width>]>>
\end{lstlisting}

\subsection{Parsing}
\label{sec:prompt-parsing}
\begin{lstlisting}[style=promptstyle]
You are an expert at analyzing early learning standards documents. You will be given a list of detected structural elements from a US-KY (2021) standards document. Each element has a level (domain, strand, sub_strand, or indicator), a code, a title, a description, and source information.

Your task is to resolve the hierarchy: assign each indicator to its correct domain, strand, and sub_strand based on the document's structural context and coding scheme.

Here are the detected elements:

[]

Return a JSON array where each object represents one indicator with its full hierarchy. Use this exact schema for each object:

{
  "domain_code": "string",
  "domain_name": "string",
  "domain_description": "string or null",
  "strand_code": "string or null",
  "strand_name": "string or null",
  "strand_description": "string or null",
  "sub_strand_code": "string or null",
  "sub_strand_name": "string or null",
  "sub_strand_description": "string or null",
  "indicator_code": "string",
  "indicator_name": "string",
  "indicator_description": "string or null",
  "age_band": "string or null",
  "column_label": "string or null",
  "source_page": integer,
  "source_text": "string"
}

Rules:
- Populate domain_description, strand_description, and sub_strand_description from the document text (the description field of the corresponding element). Use null if no description exists for that level.
- ABSENCE IS `null`, NEVER `""`. Any description you cannot fill — because the level has no description, or because the corresponding element's own description is null or blank — is `null`. Never write an empty string or a string of spaces for a description field. `""` and `null` are two spellings of the same fact, and emitting both across one document makes the same absence irreconcilable downstream. This applies to every description field in the schema.
- If a hierarchy level does not exist (e.g. no sub_strand), set its code, name, and description to null.
- For indicator_name: use the actual title of the indicator (e.g. "Curiosity and Interest"), NOT age-band/column labels like "Early", "Later", "Discovering", "PK3", "By 36 months", etc. Strip any such pre-text from the title.
- For indicator_description: use the full descriptive text of the indicator EXACTLY as it appears in the source, INCLUDING any leading proficiency label such as "Discovering:"/"Developing:"/"Broadening:" — that label carries the column's distinguishing content and MUST be kept in the description. (Age-column rows like Early/Later have no such inline label, so nothing is added.) This may be null if no description exists beyond the title.
- For age_band: examine each indicator's code, title, description, source_text, and its detected age_band field for age information. Normalize a real age RANGE to BARE months like "36-48" (PK3 → 36-48, PK4 → 48-60, "3 to 4 ½ Years" → 36-54, "4 to 5 ½ Years" → 48-66). If the column is NOT an age range — e.g. a proficiency level such as "Discovering"/"Developing"/"Broadening" — set age_band to null. The caller applies the default age band "36-60" for nulls.
- For column_label: if the indicator came from a side-by-side column, copy that column's label VERBATIM from the element's detected age_band field (e.g. "Early (3 to 4 ½ Years)", "Later (4 to 5 ½ Years)", "PK3", "Discovering"); otherwise null.
- For code: output the BASE FULL CUMULATIVE hierarchical code for every level — each child's code is its parent's code followed by the child's own segment, NOT just the final segment. A foundation with local code "1.2" under domain "AL" / strand "1.0" / sub_strand "INIT" → indicator_code "AL.1.0.INIT.1.2", sub_strand_code "AL.1.0.INIT", strand_code "AL.1.0" — never a bare "1.2" or "1.0". When a code is already fully qualified (e.g. an indicator detected as "SED.1.1.a"), use it as-is and derive the parents (strand_code "SED.1", sub_strand_code "SED.1.1").
- ALREADY-QUALIFIED codes are used AS-IS and are NEVER re-prefixed — at EVERY level, not just the indicator. An element's detected `code` is already qualified when it is a dotted path whose FIRST segment is its own domain's code (after stripping any leading column token, below): an indicator detected as "SED.1.1.a" under domain "SED", or a sub_strand detected as "AB.CD" under domain "AB". For such an element, building the cumulative chain means CONFIRMING that prefix, not prepending it a second time: a sub_strand detected "AB.CD" under domain "AB" and strand "AB.2" → sub_strand_code "AB.CD" — NEVER "AB.2.AB.CD", and NEVER "AB.2.CD". Only an element whose code is a BARE LOCAL segment ("1.2", "A", "INIT", "Concept 1") gets its parents' code prepended.
- A document's printed code NAMESPACE may SKIP a level that its hierarchy HAS. When you derive parent codes by peeling segments off a fully-qualified code, peel only as far as the namespace actually reaches: if peeling one more segment would give a level the SAME code as its own parent, that level is not in the namespace, and you must build it from its OWN heading's identifier appended to its parent's code instead. Worked example — indicators "AB.CD.PK1"/"AB.CD.PK2", a sub_strand detected "AB.CD", a strand whose detected code is "<Something> Standard 2", a domain "AB": peeling gives sub_strand_code "AB.CD" (right), but peeling again would give the strand "AB", which is already the domain's code — so the strand instead takes its heading's bare identifier "2" → strand_code "AB.2". The levels the namespace DOES cover keep the document's spelling exactly: sub_strand_code "AB.CD", indicator_code "AB.CD.PK1". The sub_strand's code then does not literally extend the strand's, and that is the correct answer: the identifier the document PRINTS is the one the standard is cited by, and it outranks cosmetic continuity of the chain.
- When an element's `code` is itself a structural label + identifier (e.g. "Strand 1", "Concept 1", "Goal 2", "Pillar A", "Unit 3" — any structural word the document uses, followed by a number or letter), use ONLY the bare identifier as that element's segment in the cumulative chain: "Strand 1" → segment "1", "Concept 1" → segment "1", "Pillar A" → segment "A". The label word merely names the level and is already captured by the element's `level`; it must NOT appear inside the cumulative `code`. Example: a strand with code "Strand 1" under domain "SED" → strand_code "SED.1" (not "SED.Strand 1"); a sub_strand with code "Concept 1" under that strand → sub_strand_code "SED.1.1" (not "SED.Strand 1.Concept 1"). Apply this to ANY label word, not just the examples.
- PRESERVE the bare identifier VERBATIM — do NOT renumber, pad, or drop any part of it. In particular, keep a decimal/dotted identifier exactly as written, INCLUDING a trailing ".0": "Strand: 1.0" → segment "1.0" (NEVER "1"), "Strand 2.0" → segment "2.0". So a strand labeled "1.0" under domain "AL" → strand_code "AL.1.0" (never "AL.1"). A trailing ".0" is part of the document's id, not a droppable minor version. (This differs from a strand whose id genuinely IS a bare integer — e.g. detected "Strand 1" or derived from an indicator code like "SED.1.1.a" → strand segment "1"; preserve whatever the id actually is.)
- A sub_strand and its child indicator must NEVER share the same code. Some documents number a named sub_strand (e.g. a "Foundation") with the SAME local number that its single child indicator also carries — e.g. a sub_strand titled "Initiative" with local code "1.2" sitting directly above an indicator also coded "1.2". When a sub_strand's local code would otherwise be identical to one of its child indicators' local codes, that shared number belongs to the INDICATOR; derive the sub_strand's OWN segment from its TITLE instead, as a ≤5-char uppercase abbreviation, using the SAME procedure the detector uses: split the title on spaces and slashes (a hyphenated compound is ONE word), DROP every connector word (`a an the and or but nor of to in on at by for from with into about over under through as`, and `&`), then — if exactly one content word remains take its first 4 letters ("Initiative" → "INIT", "Vocabulary" → "VOCA"), otherwise take the first letter of each remaining content word capped at 5 ("Concepts About Print" → "CP", "Approaches to Learning" → "AL"). Build the cumulative chain with that title-derived segment: sub_strand "Initiative" under domain "AL" / strand "1.0" → sub_strand_code "AL.1.0.INIT", and its child indicator (local code "1.2") → indicator_code "AL.1.0.INIT.1.2". A sub_strand whose code is already distinct from its indicators' (a letter, or an already-abbreviated token like "VOCA") keeps that code unchanged. This title-derived segment is a LAST RESORT, exactly as in the detector: it applies only where the document leaves the sub_strand uncoded. A sub_strand that arrives with a printed dotted code ("AB.CD") keeps it and is used as-is per the already-qualified rule above — never replaced by an abbreviation of its title.
- STRIP any leading age/column token from every indicator code: when an indicator appears in multiple side-by-side columns, each variant's detected code may begin with a token identifying its column (e.g. a grade-band prefix like `PK3.` or `PK4.`, an age-group label, or any other column-identifying token prepended to the hierarchical sequence). Strip that leading token and output only the base code shared across all column variants. Then use that stripped indicator code to derive ALL parent codes in the cumulative chain (domain, strand, sub_strand) — the stripped indicator prefix is the ground truth for the parent hierarchy, even if a detected parent element carries a different label. Example: `PK3.I.A.2` and `PK4.I.A.2` → base code `I.A.2`; domain_code=`I`, strand_code=`I.A`.
- SEPARATE domains may legitimately SHARE a strand or sub_strand TITLE, yet they remain DISTINCT entities under DISTINCT domains. For example a document can contain both a "Foundational Language Development" (FLD) domain and an "English Language Development" (ELD) domain, and BOTH may have a "Listening and Speaking" strand AND a "Vocabulary" / "Grammar" / "Phonological Awareness" sub_strand. Every code in an indicator's chain — domain_code, strand_code, AND sub_strand_code — MUST begin with the SAME domain prefix as that indicator's OWN indicator_code. NEVER borrow another domain's prefix for a strand or sub_strand just because the title matches: an indicator coded `FLD.1.0.VOCA.1.1` has strand_code `FLD.1.0` and sub_strand_code `FLD.1.0.VOCA` — NOT `ELD.1.0`/`ELD.1.0.VOCA`, even though the ELD domain has an identically-titled "Vocabulary" sub_strand. Resolve each indicator's parents strictly within its own domain.
- DISAMBIGUATE side-by-side columns by APPENDING a token to the END of the indicator code, so two variants of the same outcome (which share an identical base code) get DISTINCT indicator codes — and therefore distinct standard_ids. Append the token to the INDICATOR code ONLY; NEVER add it to the domain, strand, or sub_strand code. Choose the token by the column's type:
  - AGE-RANGE column (the cell carries an age range — e.g. "Early (3 to 4 ½ Years)", "Later (4 to 5 ½ Years)", "PK3", "PK4"): append the SAME normalized month range you put in this indicator's `age_band`, written exactly as "{start}-{end}" (no spaces, no "months"). Examples: a PK3 outcome with base `I.A.2` → `I.A.2.36-48`, its PK4 variant → `I.A.2.48-60`; an "Early" outcome with base `AL.1.0.INIT.1.2` → `AL.1.0.INIT.1.2.36-54`, its "Later" variant → `AL.1.0.INIT.1.2.48-66`. Apply this even when only one age column is present for the outcome (a lone PK4 outcome `VI.A.1` still becomes `VI.A.1.48-60`).
  - NON-AGE column (a proficiency or similar label where the columns share one age range, so `age_band` is null — e.g. "Discovering"/"Developing"/"Broadening"): append the FIRST FOUR LETTERS of the column label, UPPERCASED. Examples: base `ELD.1.0.VOCA.1.1` → `ELD.1.0.VOCA.1.1.DISC` (Discovering), `ELD.1.0.VOCA.1.1.DEVE` (Developing), `ELD.1.0.VOCA.1.1.BROA` (Broadening).
  - If the indicator does NOT come from an age/column cell (no per-column `age_band` and no `column_label`), append NOTHING — leave the base code as-is.
- Return ONLY the JSON array, no other text.
- Every indicator element must appear exactly once in the output.
- There will be cases where you see "No PK3 outcomes for this domain of learning." or similar wording. This means that for the given indicator and age, there is no outcome. In this case, the indicator should be omitted. Do not try to attach it to another indicator.

CRITICAL — RESOLVING HIERARCHY USING CODES AND CONTEXT:
The detected elements come from processing the document in overlapping chunks. This means:
- A strand or sub_strand may have been detected in one chunk while its child indicators were detected in a different chunk. You MUST still link them correctly using code prefixes, naming patterns, and document order.
- Use the hierarchical code structure to resolve parentage. For example, indicators with codes like "PK3.I.A.1" and "PK4.I.A.1" belong under the strand with code "A" (Self-Concept) in domain "I" (Social and Emotional Development).
- If a strand or sub_strand appears in the elements list but has no indicators directly following it, look for indicators elsewhere in the list whose codes or source context indicate they belong under that strand/sub_strand.
- Do NOT drop strands or sub_strands just because their indicators are not adjacent in the element list. Every strand and sub_strand should appear as the parent of at least one indicator in the output (unless it genuinely has no indicators in the document).
- ASSIGN every indicator to the most-recent preceding sub_strand in document order within its strand. A sub_strand header's coverage extends FORWARD across ALL following indicators until the next sub_strand header, the next strand, or the next domain begins. Do this even when those later indicators' local numbers do not look "under" the sub_strand's own number — documents commonly number sub_strands and indicators in ONE flat sequence, so a "Grammar" sub_strand introduced at 1.4 covers foundations 1.4, 1.5, 1.6, 1.7, AND 1.8 right up until "Strand 2.0" begins. NEVER leave an indicator's sub_strand null when a sub_strand precedes it within the same strand; only a strand whose indicators appear before ANY sub_strand header (or a document with no sub_strands at all) yields a null sub_strand.
\end{lstlisting}
